\chapter{实验性功能}

% === 源文件: experiments/onboarding-config-protocol.md ===
\section{新手引导 + 配置协议}


目的：CLI、macOS 应用和 Web UI 之间共享的新手引导 + 配置界面。

\subsection{组件}


\begin{itemize}
  \item 向导引擎（共享会话 + 提示 + 新手引导状态）。
  \item CLI 新手引导使用与 UI 客户端相同的向导流程。
  \item Gateway 网关 RPC 公开向导 + 配置模式端点。
  \item macOS 新手引导使用向导步骤模型。
  \item Web UI 从 JSON Schema + UI 提示渲染配置表单。
\end{itemize}


\subsection{Gateway 网关 RPC}


\begin{itemize}
  \item \texttt{wizard.start} 参数：\texttt{\{ mode?: "local"|"remote", workspace?: string \}}
  \item \texttt{wizard.next} 参数：\texttt{\{ sessionId, answer?: \{ stepId, value? \} \}}
  \item \texttt{wizard.cancel} 参数：\texttt{\{ sessionId \}}
  \item \texttt{wizard.status} 参数：\texttt{\{ sessionId \}}
  \item \texttt{config.schema} 参数：\texttt{\{\}}
\end{itemize}


响应（结构）

\begin{itemize}
  \item 向导：\texttt{\{ sessionId, done, step?, status?, error? \}}
  \item 配置模式：\texttt{\{ schema, uiHints, version, generatedAt \}}
\end{itemize}


\subsection{UI 提示}


\begin{itemize}
  \item \texttt{uiHints} 按路径键入；可选元数据（label/help/group/order/advanced/sensitive/placeholder）。
  \item 敏感字段渲染为密码输入；无脱敏层。
  \item 不支持的模式节点回退到原始 JSON 编辑器。
\end{itemize}


\subsection{注意}


\begin{itemize}
  \item 本文档是跟踪新手引导/配置协议重构的唯一位置。
\end{itemize}



\clearpage

% === 源文件: experiments/plans/cron-add-hardening.md ===
\section{Cron Add 加固 \& Schema 对齐}


\subsection{背景}


最近的 Gateway 网关日志显示重复的 \texttt{cron.add} 失败，参数无效（缺少 \texttt{sessionTarget}、\texttt{wakeMode}、\texttt{payload}，以及格式错误的 \texttt{schedule}）。这表明至少有一个客户端（可能是智能体工具调用路径）正在发送包装的或部分指定的任务负载。另外，TypeScript 中的 cron 提供商枚举、Gateway 网关 schema、CLI 标志和 UI 表单类型之间存在漂移，加上 \texttt{cron.status} 的 UI 不匹配（期望 \texttt{jobCount} 而 Gateway 网关返回 \texttt{jobs}）。

\subsection{目标}


\begin{itemize}
  \item 通过规范化常见的包装负载并推断缺失的 \texttt{kind} 字段来停止 \texttt{cron.add} INVALID\_REQUEST 垃圾。
  \item 在 Gateway 网关 schema、cron 类型、CLI 文档和 UI 表单之间对齐 cron 提供商列表。
  \item 使智能体 cron 工具 schema 明确，以便 LLM 生成正确的任务负载。
  \item 修复 Control UI cron 状态任务计数显示。
  \item 添加测试以覆盖规范化和工具行为。
\end{itemize}


\subsection{非目标}


\begin{itemize}
  \item 更改 cron 调度语义或任务执行行为。
  \item 添加新的调度类型或 cron 表达式解析。
  \item 除了必要的字段修复外，不大改 cron 的 UI/UX。
\end{itemize}


\subsection{发现（当前差距）}


\begin{itemize}
  \item Gateway 网关中的 \texttt{CronPayloadSchema} 排除了 \texttt{signal} + \texttt{imessage}，而 TS 类型包含它们。
  \item Control UI CronStatus 期望 \texttt{jobCount}，但 Gateway 网关返回 \texttt{jobs}。
  \item 智能体 cron 工具 schema 允许任意 \texttt{job} 对象，导致格式错误的输入。
  \item Gateway 网关严格验证 \texttt{cron.add} 而不进行规范化，因此包装的负载会失败。
\end{itemize}


\subsection{变更内容}


\begin{itemize}
  \item \texttt{cron.add} 和 \texttt{cron.update} 现在规范化常见的包装形式并推断缺失的 \texttt{kind} 字段。
  \item 智能体 cron 工具 schema 与 Gateway 网关 schema 匹配，减少无效负载。
  \item 提供商枚举在 Gateway 网关、CLI、UI 和 macOS 选择器之间对齐。
  \item Control UI 使用 Gateway 网关的 \texttt{jobs} 计数字段显示状态。
\end{itemize}


\subsection{当前行为}


\begin{itemize}
  \item \textbf{规范化：}包装的 \texttt{data}/\texttt{job} 负载被解包；\texttt{schedule.kind} 和 \texttt{payload.kind} 在安全时被推断。
  \item \textbf{默认值：}当缺失时，为 \texttt{wakeMode} 和 \texttt{sessionTarget} 应用安全默认值。
  \item \textbf{提供商：}Discord/Slack/Signal/iMessage 现在在 CLI/UI 中一致显示。
\end{itemize}


参见 \href{/automation/cron-jobs}{Cron 任务} 了解规范化的形式和示例。

\subsection{验证}


\begin{itemize}
  \item 观察 Gateway 网关日志中 \texttt{cron.add} INVALID\_REQUEST 错误是否减少。
  \item 确认 Control UI cron 状态在刷新后显示任务计数。
\end{itemize}


\subsection{可选后续工作}


\begin{itemize}
  \item 手动 Control UI 冒烟测试：为每个提供商添加一个 cron 任务 + 验证状态任务计数。
\end{itemize}


\subsection{开放问题}


\begin{itemize}
  \item \texttt{cron.add} 是否应该接受来自客户端的显式 \texttt{state}（当前被 schema 禁止）？
  \item 我们是否应该允许 \texttt{webchat} 作为显式投递提供商（当前在投递解析中被过滤）？
\end{itemize}



\clearpage

% === 源文件: experiments/plans/group-policy-hardening.md ===
\section{Telegram 允许列表加固}


\textbf{日期}：2026-01-05
\textbf{状态}：已完成
\textbf{PR}：\#216

\subsection{摘要}


Telegram 允许列表现在不区分大小写地接受 \texttt{telegram:} 和 \texttt{tg:} 前缀，并容忍意外的空白。这使入站允许列表检查与出站发送规范化保持一致。

\subsection{更改内容}


\begin{itemize}
  \item 前缀 \texttt{telegram:} 和 \texttt{tg:} 被同等对待（不区分大小写）。
  \item 允许列表条目会被修剪；空条目会被忽略。
\end{itemize}


\subsection{示例}


以下所有形式都被接受为同一 ID：

\begin{itemize}
  \item \texttt{telegram:123456}
  \item \texttt{TG:123456}
  \item \texttt{tg:123456}
\end{itemize}


\subsection{为什么重要}


从日志或聊天 ID 复制/粘贴通常会包含前缀和空白。规范化可避免在决定是否在私信或群组中响应时出现误判。

\subsection{相关文档}


\begin{itemize}
  \item \href{/channels/groups}{群聊}
  \item \href{/channels/telegram}{Telegram 提供商}
\end{itemize}



\clearpage

% === 源文件: experiments/plans/openresponses-gateway.md ===
\section{OpenResponses Gateway 网关集成计划}


\subsection{背景}


OpenClaw Gateway 网关目前在 \texttt{/v1/chat/completions} 暴露了一个最小的 OpenAI 兼容 Chat Completions 端点（参见 \href{/gateway/openai-http-api}{OpenAI Chat Completions}）。

Open Responses 是基于 OpenAI Responses API 的开放推理标准。它专为智能体工作流设计，使用基于项目的输入加语义流式事件。OpenResponses 规范定义的是 \texttt{/v1/responses}，而不是 \texttt{/v1/chat/completions}。

\subsection{目标}


\begin{itemize}
  \item 添加一个遵循 OpenResponses 语义的 \texttt{/v1/responses} 端点。
  \item 保留 Chat Completions 作为兼容层，易于禁用并最终移除。
  \item 使用隔离的、可复用的 schema 标准化验证和解析。
\end{itemize}


\subsection{非目标}


\begin{itemize}
  \item 第一阶段完全实现 OpenResponses 功能（图片、文件、托管工具）。
  \item 替换内部智能体执行逻辑或工具编排。
  \item 在第一阶段更改现有的 \texttt{/v1/chat/completions} 行为。
\end{itemize}


\subsection{研究摘要}


来源：OpenResponses OpenAPI、OpenResponses 规范网站和 Hugging Face 博客文章。

提取的关键点：

\begin{itemize}
  \item \texttt{POST /v1/responses} 接受 \texttt{CreateResponseBody} 字段，如 \texttt{model}、\texttt{input}（字符串或 \texttt{ItemParam[]}）、\texttt{instructions}、\texttt{tools}、\texttt{tool\_choice}、\texttt{stream}、\texttt{max\_output\_tokens} 和 \texttt{max\_tool\_calls}。
  \item \texttt{ItemParam} 是以下类型的可区分联合：
  \item 具有角色 \texttt{system}、\texttt{developer}、\texttt{user}、\texttt{assistant} 的 \texttt{message} 项
  \item \texttt{function\_call} 和 \texttt{function\_call\_output}
  \item \texttt{reasoning}
  \item \texttt{item\_reference}
  \item 成功响应返回带有 \texttt{object: "response"}、\texttt{status} 和 \texttt{output} 项的 \texttt{ResponseResource}。
  \item 流式传输使用语义事件，如：
  \item \texttt{response.created}、\texttt{response.in\_progress}、\texttt{response.completed}、\texttt{response.failed}
  \item \texttt{response.output\_item.added}、\texttt{response.output\_item.done}
  \item \texttt{response.content\_part.added}、\texttt{response.content\_part.done}
  \item \texttt{response.output\_text.delta}、\texttt{response.output\_text.done}
  \item 规范要求：
  \item \texttt{Content-Type: text/event-stream}
  \item \texttt{event:} 必须匹配 JSON \texttt{type} 字段
  \item 终止事件必须是字面量 \texttt{[DONE]}
  \item Reasoning 项可能暴露 \texttt{content}、\texttt{encrypted\_content} 和 \texttt{summary}。
  \item HF 示例在请求中包含 \texttt{OpenResponses-Version: latest}（可选头部）。
\end{itemize}


\subsection{提议的架构}


\begin{itemize}
  \item 添加 \texttt{src/gateway/open-responses.schema.ts}，仅包含 Zod schema（无 gateway 导入）。
  \item 添加 \texttt{src/gateway/openresponses-http.ts}（或 \texttt{open-responses-http.ts}）用于 \texttt{/v1/responses}。
  \item 保持 \texttt{src/gateway/openai-http.ts} 不变，作为遗留兼容适配器。
  \item 添加配置 \texttt{gateway.http.endpoints.responses.enabled}（默认 \texttt{false}）。
  \item 保持 \texttt{gateway.http.endpoints.chatCompletions.enabled} 独立；允许两个端点分别切换。
  \item 当 Chat Completions 启用时发出启动警告，以表明其遗留状态。
\end{itemize}


\subsection{Chat Completions 弃用路径}


\begin{itemize}
  \item 保持严格的模块边界：responses 和 chat completions 之间不共享 schema 类型。
  \item 通过配置使 Chat Completions 成为可选，这样无需代码更改即可禁用。
  \item 一旦 \texttt{/v1/responses} 稳定，更新文档将 Chat Completions 标记为遗留。
  \item 可选的未来步骤：将 Chat Completions 请求映射到 Responses 处理器，以便更简单地移除。
\end{itemize}


\subsection{第一阶段支持子集}


\begin{itemize}
  \item 接受 \texttt{input} 为字符串或带有消息角色和 \texttt{function\_call\_output} 的 \texttt{ItemParam[]}。
  \item 将 system 和 developer 消息提取到 \texttt{extraSystemPrompt} 中。
  \item 使用最近的 \texttt{user} 或 \texttt{function\_call\_output} 作为智能体运行的当前消息。
  \item 对不支持的内容部分（图片/文件）返回 \texttt{invalid\_request\_error} 拒绝。
  \item 返回带有 \texttt{output\_text} 内容的单个助手消息。
  \item 返回带有零值的 \texttt{usage}，直到 token 计数接入。
\end{itemize}


\subsection{验证策略（无 SDK）}


\begin{itemize}
  \item 为以下支持子集实现 Zod schema：
  \item \texttt{CreateResponseBody}
  \item \texttt{ItemParam} + 消息内容部分联合
  \item \texttt{ResponseResource}
  \item Gateway 网关使用的流式事件形状
  \item 将 schema 保存在单个隔离模块中，以避免漂移并允许未来代码生成。
\end{itemize}


\subsection{流式实现（第一阶段）}


\begin{itemize}
  \item 带有 \texttt{event:} 和 \texttt{data:} 的 SSE 行。
  \item 所需序列（最小可行）：
  \item \texttt{response.created}
  \item \texttt{response.output\_item.added}
  \item \texttt{response.content\_part.added}
  \item \texttt{response.output\_text.delta}（根据需要重复）
  \item \texttt{response.output\_text.done}
  \item \texttt{response.content\_part.done}
  \item \texttt{response.completed}
  \item \texttt{[DONE]}
\end{itemize}


\subsection{测试和验证计划}


\begin{itemize}
  \item 为 \texttt{/v1/responses} 添加端到端覆盖：
  \item 需要认证
  \item 非流式响应形状
  \item 流式事件顺序和 \texttt{[DONE]}
  \item 使用头部和 \texttt{user} 的会话路由
  \item 保持 \texttt{src/gateway/openai-http.e2e.test.ts} 不变。
  \item 手动：用 \texttt{stream: true} curl \texttt{/v1/responses} 并验证事件顺序和终止 \texttt{[DONE]}。
\end{itemize}


\subsection{文档更新（后续）}


\begin{itemize}
  \item 为 \texttt{/v1/responses} 使用和示例添加新文档页面。
  \item 更新 \texttt{/gateway/openai-http-api}，添加遗留说明和指向 \texttt{/v1/responses} 的指针。
\end{itemize}



\clearpage

% === 源文件: experiments/proposals/model-config.md ===
\section{模型配置（探索）}


本文档记录了未来模型配置的\textbf{构想}。这不是正式的发布规范。如需了解当前行为，请参阅：

\begin{itemize}
  \item \href{/concepts/models}{模型}
  \item \href{/concepts/model-failover}{模型故障转移}
  \item \href{/concepts/oauth}{OAuth + 配置文件}
\end{itemize}


\subsection{动机}


运营者希望：

\begin{itemize}
  \item 每个提供商支持多个认证配置文件（个人 vs 工作）。
  \item 简单的 \texttt{/model} 选择，并具有可预测的回退行为。
  \item 文本模型与图像模型之间有清晰的分离。
\end{itemize}


\subsection{可能的方向（高层级）}


\begin{itemize}
  \item 保持模型选择简洁：\texttt{provider/model} 加可选别名。
  \item 允许提供商拥有多个认证配置文件，并指定明确的顺序。
  \item 使用全局回退列表，使所有会话以一致的方式进行故障转移。
  \item 仅在明确配置时才覆盖图像路由。
\end{itemize}


\subsection{待解决的问题}


\begin{itemize}
  \item 配置文件轮换应该按提供商还是按模型进行？
  \item UI 应如何为会话展示配置文件选择？
  \item 从旧版配置键迁移的最安全路径是什么？
\end{itemize}



\clearpage

% === 源文件: experiments/research/memory.md ===
\section{工作区记忆 v2（离线）：研究笔记}


目标：Clawd 风格的工作区（\texttt{agents.defaults.workspace}，默认 \texttt{\textasciitilde{}/.openclaw/workspace}），其中"记忆"以每天一个 Markdown 文件（\texttt{memory/YYYY-MM-DD.md}）加上一小组稳定文件（例如 \texttt{memory.md}、\texttt{SOUL.md}）的形式存储。

本文档提出一种\textbf{离线优先}的记忆架构，保持 Markdown 作为规范的、可审查的数据源，但通过派生索引添加\textbf{结构化回忆}（搜索、实体摘要、置信度更新）。

\subsection{为什么要改变？}


当前设置（每天一个文件）非常适合：

\begin{itemize}
  \item "仅追加"式日志记录
  \item 人工编辑
  \item git 支持的持久性 + 可审计性
  \item 低摩擦捕获（"直接写下来"）
\end{itemize}


但它在以下方面较弱：

\begin{itemize}
  \item 高召回率检索（"我们对 X 做了什么决定？"、"上次我们尝试 Y 时？"）
  \item 以实体为中心的答案（"告诉我关于 Alice / The Castle / warelay 的信息"）而无需重读多个文件
  \item 观点/偏好稳定性（以及变化时的证据）
  \item 时间约束（"2025 年 11 月期间什么是真实的？"）和冲突解决
\end{itemize}


\subsection{设计目标}


\begin{itemize}
  \item \textbf{离线}：无需网络即可工作；可在笔记本电脑/Castle 上运行；无云依赖。
  \item \textbf{可解释}：检索的项目应该可归因（文件 + 位置）并与推理分离。
  \item \textbf{低仪式感}：每日日志保持 Markdown，无需繁重的 schema 工作。
  \item \textbf{增量式}：v1 仅使用 FTS 就很有用；语义/向量和图是可选升级。
  \item \textbf{对智能体友好}：使"在 token 预算内回忆"变得简单（返回小型事实包）。
\end{itemize}


\subsection{北极星模型（Hindsight × Letta）}


需要融合两个部分：

\begin{enumerate}
  \item \textbf{Letta/MemGPT 风格的控制循环}
\end{enumerate}


\begin{itemize}
  \item 保持一个小的"核心"始终在上下文中（角色 + 关键用户事实）
  \item 其他所有内容都在上下文之外，通过工具检索
  \item 记忆写入是显式的工具调用（append/replace/insert），持久化后在下一轮重新注入
\end{itemize}


\begin{enumerate}
  \item \textbf{Hindsight 风格的记忆基底}
\end{enumerate}


\begin{itemize}
  \item 分离观察到的、相信的和总结的内容
  \item 支持 retain/recall/reflect
  \item 带有置信度的观点可以随证据演变
  \item 实体感知检索 + 时间查询（即使没有完整的知识图谱）
\end{itemize}


\subsection{提议的架构（Markdown 数据源 + 派生索引）}


\subsubsection{规范存储（git 友好）}


保持 \texttt{\textasciitilde{}/.openclaw/workspace} 作为规范的人类可读记忆。

建议的工作区布局：

\begin{lstlisting}
~/.openclaw/workspace/
  memory.md                    # 小型：持久事实 + 偏好（类似核心）
  memory/
    YYYY-MM-DD.md              # 每日日志（追加；叙事）
  bank/                        # "类型化"记忆页面（稳定、可审查）
    world.md                   # 关于世界的客观事实
    experience.md              # 智能体做了什么（第一人称）
    opinions.md                # 主观偏好/判断 + 置信度 + 证据指针
    entities/
      Peter.md
      The-Castle.md
      warelay.md
      ...
\end{lstlisting}


注意：

\begin{itemize}
  \item \textbf{每日日志保持为每日日志}。无需将其转换为 JSON。
  \item \texttt{bank/} 文件是\textbf{经过整理的}，由反思任务生成，仍可手动编辑。
  \item \texttt{memory.md} 保持"小型 + 类似核心"：你希望 Clawd 每次会话都能看到的内容。
\end{itemize}


\subsubsection{派生存储（机器回忆）}


在工作区下添加派生索引（不一定需要 git 跟踪）：

\begin{lstlisting}
~/.openclaw/workspace/.memory/index.sqlite
\end{lstlisting}


后端支持：

\begin{itemize}
  \item 用于事实 + 实体链接 + 观点元数据的 SQLite schema
  \item SQLite \textbf{FTS5} 用于词法回忆（快速、小巧、离线）
  \item 可选的嵌入表用于语义回忆（仍然离线）
\end{itemize}


索引始终\textbf{可从 Markdown 重建}。

\subsection{Retain / Recall / Reflect（操作循环）}


\subsubsection{Retain：将每日日志规范化为"事实"}


Hindsight 在这里重要的关键洞察：存储\textbf{叙事性、自包含的事实}，而不是微小的片段。

\texttt{memory/YYYY-MM-DD.md} 的实用规则：

\begin{itemize}
  \item 在一天结束时（或期间），添加一个 \texttt{\#\# Retain} 部分，包含 2-5 个要点：
  \item 叙事性（保留跨轮上下文）
  \item 自包含（独立时也有意义）
  \item 标记类型 + 实体提及
\end{itemize}


示例：

\begin{lstlisting}
## Retain
- W @Peter: Currently in Marrakech (Nov 27–Dec 1, 2025) for Andy's birthday.
- B @warelay: I fixed the Baileys WS crash by wrapping connection.update handlers in try/catch (see memory/2025-11-27.md).
- O(c=0.95) @Peter: Prefers concise replies (&lt;1500 chars) on WhatsApp; long content goes into files.
\end{lstlisting}


最小化解析：

\begin{itemize}
  \item 类型前缀：\texttt{W}（世界）、\texttt{B}（经历/传记）、\texttt{O}（观点）、\texttt{S}（观察/摘要；通常是生成的）
  \item 实体：\texttt{@Peter}、\texttt{@warelay} 等（slug 映射到 \texttt{bank/entities/*.md}）
  \item 观点置信度：\texttt{O(c=0.0..1.0)} 可选
\end{itemize}


如果你不想让作者考虑这些：反思任务可以从日志的其余部分推断这些要点，但有一个显式的 \texttt{\#\# Retain} 部分是最简单的"质量杠杆"。

\subsubsection{Recall：对派生索引的查询}


Recall 应支持：

\begin{itemize}
  \item \textbf{词法}："查找精确的术语/名称/命令"（FTS5）
  \item \textbf{实体}："告诉我关于 X 的信息"（实体页面 + 实体链接的事实）
  \item \textbf{时间}："11 月 27 日前后发生了什么"/"自上周以来"
  \item \textbf{观点}："Peter 偏好什么？"（带置信度 + 证据）
\end{itemize}


返回格式应对智能体友好并引用来源：

\begin{itemize}
  \item \texttt{kind}（\texttt{world|experience|opinion|observation}）
  \item \texttt{timestamp}（来源日期，或如果存在则提取的时间范围）
  \item \texttt{entities}（\texttt{["Peter","warelay"]}）
  \item \texttt{content}（叙事性事实）
  \item \texttt{source}（\texttt{memory/2025-11-27.md\#L12} 等）
\end{itemize}


\subsubsection{Reflect：生成稳定页面 + 更新信念}


反思是一个定时任务（每日或心跳 \texttt{ultrathink}），它：

\begin{itemize}
  \item 根据最近的事实更新 \texttt{bank/entities/*.md}（实体摘要）
  \item 根据强化/矛盾更新 \texttt{bank/opinions.md} 置信度
  \item 可选地提议对 \texttt{memory.md}（"类似核心"的持久事实）的编辑
\end{itemize}


观点演变（简单、可解释）：

\begin{itemize}
  \item 每个观点有：
  \item 陈述
  \item 置信度 \texttt{c ∈ [0,1]}
  \item last\_updated
  \item 证据链接（支持 + 矛盾的事实 ID）
  \item 当新事实到达时：
  \item 通过实体重叠 + 相似性找到候选观点（先 FTS，后嵌入）
  \item 通过小幅增量更新置信度；大幅跳跃需要强矛盾 + 重复证据
\end{itemize}


\subsection{CLI 集成：独立 vs 深度集成}


建议：\textbf{深度集成到 OpenClaw}，但保持可分离的核心库。

\subsubsection{为什么要集成到 OpenClaw？}


\begin{itemize}
  \item OpenClaw 已经知道：
  \item 工作区路径（\texttt{agents.defaults.workspace}）
  \item 会话模型 + 心跳
  \item 日志记录 + 故障排除模式
  \item 你希望智能体自己调用工具：
  \item \texttt{openclaw memory recall "…" --k 25 --since 30d}
  \item \texttt{openclaw memory reflect --since 7d}
\end{itemize}


\subsubsection{为什么仍要分离库？}


\begin{itemize}
  \item 保持记忆逻辑可测试，无需 Gateway 网关/运行时
  \item 可从其他上下文重用（本地脚本、未来的桌面应用等）
\end{itemize}


形态：
记忆工具预计是一个小型 CLI + 库层，但这仅是探索性的。

\subsection{"S-Collide" / SuCo：何时使用（研究）}


如果"S-Collide"指的是 \textbf{SuCo（Subspace Collision）}：这是一种 ANN 检索方法，通过在子空间中使用学习/结构化碰撞来实现强召回/延迟权衡（论文：arXiv 2411.14754，2024）。

对于 \texttt{\textasciitilde{}/.openclaw/workspace} 的务实观点：

\begin{itemize}
  \item \textbf{不要从} SuCo 开始。
  \item 从 SQLite FTS +（可选的）简单嵌入开始；你会立即获得大部分 UX 收益。
  \item 仅在以下情况下考虑 SuCo/HNSW/ScaNN 级别的解决方案：
  \item 语料库很大（数万/数十万个块）
  \item 暴力嵌入搜索变得太慢
  \item 召回质量明显受到词法搜索的瓶颈限制
\end{itemize}


离线友好的替代方案（按复杂性递增）：

\begin{itemize}
  \item SQLite FTS5 + 元数据过滤（零 ML）
  \item 嵌入 + 暴力搜索（如果块数量低，效果出奇地好）
  \item HNSW 索引（常见、稳健；需要库绑定）
  \item SuCo（研究级；如果有可嵌入的可靠实现则很有吸引力）
\end{itemize}


开放问题：

\begin{itemize}
  \item 对于你的机器（笔记本 + 台式机）上的"个人助理记忆"，\textbf{最佳}的离线嵌入模型是什么？
  \item 如果你已经有 Ollama：使用本地模型嵌入；否则在工具链中附带一个小型嵌入模型。
\end{itemize}


\subsection{最小可用试点}


如果你想要一个最小但仍有用的版本：

\begin{itemize}
  \item 添加 \texttt{bank/} 实体页面和每日日志中的 \texttt{\#\# Retain} 部分。
  \item 使用 SQLite FTS 进行带引用的回忆（路径 + 行号）。
  \item 仅在召回质量或规模需要时添加嵌入。
\end{itemize}


\subsection{参考资料}


\begin{itemize}
  \item Letta / MemGPT 概念："核心记忆块" + "档案记忆" + 工具驱动的自编辑记忆。
  \item Hindsight 技术报告："retain / recall / reflect"，四网络记忆，叙事性事实提取，观点置信度演变。
  \item SuCo：arXiv 2411.14754（2024）："Subspace Collision"近似最近邻检索。
\end{itemize}



\clearpage
% === 源文件: experiments/plans/acp-thread-bound-agents.md ===
\section{ACP Thread Bound Agents}


\subsection{Overview}


This plan defines how OpenClaw should support ACP coding agents in thread-capable channels (Discord first) with production-level lifecycle and recovery.

Related document:

\begin{itemize}
  \item \href{/experiments/plans/acp-unified-streaming-refactor}{Unified Runtime Streaming Refactor Plan}
\end{itemize}


Target user experience:

\begin{itemize}
  \item a user spawns or focuses an ACP session into a thread
  \item user messages in that thread route to the bound ACP session
  \item agent output streams back to the same thread persona
  \item session can be persistent or one shot with explicit cleanup controls
\end{itemize}


\subsection{Decision summary}


Long term recommendation is a hybrid architecture:

\begin{itemize}
  \item OpenClaw core owns ACP control plane concerns
  \item session identity and metadata
  \item thread binding and routing decisions
  \item delivery invariants and duplicate suppression
  \item lifecycle cleanup and recovery semantics
  \item ACP runtime backend is pluggable
  \item first backend is an acpx-backed plugin service
  \item runtime does ACP transport, queueing, cancel, reconnect
\end{itemize}


OpenClaw should not reimplement ACP transport internals in core.
OpenClaw should not rely on a pure plugin-only interception path for routing.

\subsection{North-star architecture (holy grail)}


Treat ACP as a first-class control plane in OpenClaw, with pluggable runtime adapters.

Non-negotiable invariants:

\begin{itemize}
  \item every ACP thread binding references a valid ACP session record
  \item every ACP session has explicit lifecycle state (\texttt{creating}, \texttt{idle}, \texttt{running}, \texttt{cancelling}, \texttt{closed}, \texttt{error})
  \item every ACP run has explicit run state (\texttt{queued}, \texttt{running}, \texttt{completed}, \texttt{failed}, \texttt{cancelled})
  \item spawn, bind, and initial enqueue are atomic
  \item command retries are idempotent (no duplicate runs or duplicate Discord outputs)
  \item bound-thread channel output is a projection of ACP run events, never ad-hoc side effects
\end{itemize}


Long-term ownership model:

\begin{itemize}
  \item \texttt{AcpSessionManager} is the single ACP writer and orchestrator
  \item manager lives in gateway process first; can be moved to a dedicated sidecar later behind the same interface
  \item per ACP session key, manager owns one in-memory actor (serialized command execution)
  \item adapters (\texttt{acpx}, future backends) are transport/runtime implementations only
\end{itemize}


Long-term persistence model:

\begin{itemize}
  \item move ACP control-plane state to a dedicated SQLite store (WAL mode) under OpenClaw state dir
  \item keep \texttt{SessionEntry.acp} as compatibility projection during migration, not source-of-truth
  \item store ACP events append-only to support replay, crash recovery, and deterministic delivery
\end{itemize}


\subsubsection{Delivery strategy (bridge to holy-grail)}


\begin{itemize}
  \item short-term bridge
  \item keep current thread binding mechanics and existing ACP config surface
  \item fix metadata-gap bugs and route ACP turns through a single core ACP branch
  \item add idempotency keys and fail-closed routing checks immediately
  \item long-term cutover
  \item move ACP source-of-truth to control-plane DB + actors
  \item make bound-thread delivery purely event-projection based
  \item remove legacy fallback behavior that depends on opportunistic session-entry metadata
\end{itemize}


\subsection{Why not pure plugin only}


Current plugin hooks are not sufficient for end to end ACP session routing without core changes.

\begin{itemize}
  \item inbound routing from thread binding resolves to a session key in core dispatch first
  \item message hooks are fire-and-forget and cannot short-circuit the main reply path
  \item plugin commands are good for control operations but not for replacing core per-turn dispatch flow
\end{itemize}


Result:

\begin{itemize}
  \item ACP runtime can be pluginized
  \item ACP routing branch must exist in core
\end{itemize}


\subsection{Existing foundation to reuse}


Already implemented and should remain canonical:

\begin{itemize}
  \item thread binding target supports \texttt{subagent} and \texttt{acp}
  \item inbound thread routing override resolves by binding before normal dispatch
  \item outbound thread identity via webhook in reply delivery
  \item \texttt{/focus} and \texttt{/unfocus} flow with ACP target compatibility
  \item persistent binding store with restore on startup
  \item unbind lifecycle on archive, delete, unfocus, reset, and delete
\end{itemize}


This plan extends that foundation rather than replacing it.

\subsection{Architecture}


\subsubsection{Boundary model}


Core (must be in OpenClaw core):

\begin{itemize}
  \item ACP session-mode dispatch branch in the reply pipeline
  \item delivery arbitration to avoid parent plus thread duplication
  \item ACP control-plane persistence (with \texttt{SessionEntry.acp} compatibility projection during migration)
  \item lifecycle unbind and runtime detach semantics tied to session reset/delete
\end{itemize}


Plugin backend (acpx implementation):

\begin{itemize}
  \item ACP runtime worker supervision
  \item acpx process invocation and event parsing
  \item ACP command handlers (\texttt{/acp ...}) and operator UX
  \item backend-specific config defaults and diagnostics
\end{itemize}


\subsubsection{Runtime ownership model}


\begin{itemize}
  \item one gateway process owns ACP orchestration state
  \item ACP execution runs in supervised child processes via acpx backend
  \item process strategy is long lived per active ACP session key, not per message
\end{itemize}


This avoids startup cost on every prompt and keeps cancel and reconnect semantics reliable.

\subsubsection{Core runtime contract}


Add a core ACP runtime contract so routing code does not depend on CLI details and can switch backends without changing dispatch logic:

\begin{lstlisting}[language=typescript]
export type AcpRuntimePromptMode = "prompt" | "steer";

export type AcpRuntimeHandle = {
  sessionKey: string;
  backend: string;
  runtimeSessionName: string;
};

export type AcpRuntimeEvent =
  | { type: "text_delta"; stream: "output" | "thought"; text: string }
  | { type: "tool_call"; name: string; argumentsText: string }
  | { type: "done"; usage?: Record<string, number> }
  | { type: "error"; code: string; message: string; retryable?: boolean };

export interface AcpRuntime {
  ensureSession(input: {
    sessionKey: string;
    agent: string;
    mode: "persistent" | "oneshot";
    cwd?: string;
    env?: Record<string, string>;
    idempotencyKey: string;
  }): Promise<AcpRuntimeHandle>;

  submit(input: {
    handle: AcpRuntimeHandle;
    text: string;
    mode: AcpRuntimePromptMode;
    idempotencyKey: string;
  }): Promise<{ runtimeRunId: string }>;

  stream(input: {
    handle: AcpRuntimeHandle;
    runtimeRunId: string;
    onEvent: (event: AcpRuntimeEvent) => Promise<void> | void;
    signal?: AbortSignal;
  }): Promise<void>;

  cancel(input: {
    handle: AcpRuntimeHandle;
    runtimeRunId?: string;
    reason?: string;
    idempotencyKey: string;
  }): Promise<void>;

  close(input: { handle: AcpRuntimeHandle; reason: string; idempotencyKey: string }): Promise<void>;

  health?(): Promise<{ ok: boolean; details?: string }>;
}
\end{lstlisting}


Implementation detail:

\begin{itemize}
  \item first backend: \texttt{AcpxRuntime} shipped as a plugin service
  \item core resolves runtime via registry and fails with explicit operator error when no ACP runtime backend is available
\end{itemize}


\subsubsection{Control-plane data model and persistence}


Long-term source-of-truth is a dedicated ACP SQLite database (WAL mode), for transactional updates and crash-safe recovery:

\begin{itemize}
  \item \texttt{acp\_sessions}
  \item \texttt{session\_key} (pk), \texttt{backend}, \texttt{agent}, \texttt{mode}, \texttt{cwd}, \texttt{state}, \texttt{created\_at}, \texttt{updated\_at}, \texttt{last\_error}
  \item \texttt{acp\_runs}
  \item \texttt{run\_id} (pk), \texttt{session\_key} (fk), \texttt{state}, \texttt{requester\_message\_id}, \texttt{idempotency\_key}, \texttt{started\_at}, \texttt{ended\_at}, \texttt{error\_code}, \texttt{error\_message}
  \item \texttt{acp\_bindings}
  \item \texttt{binding\_key} (pk), \texttt{thread\_id}, \texttt{channel\_id}, \texttt{account\_id}, \texttt{session\_key} (fk), \texttt{expires\_at}, \texttt{bound\_at}
  \item \texttt{acp\_events}
  \item \texttt{event\_id} (pk), \texttt{run\_id} (fk), \texttt{seq}, \texttt{kind}, \texttt{payload\_json}, \texttt{created\_at}
  \item \texttt{acp\_delivery\_checkpoint}
  \item \texttt{run\_id} (pk/fk), \texttt{last\_event\_seq}, \texttt{last\_discord\_message\_id}, \texttt{updated\_at}
  \item \texttt{acp\_idempotency}
  \item \texttt{scope}, \texttt{idempotency\_key}, \texttt{result\_json}, \texttt{created\_at}, unique \texttt{(scope, idempotency\_key)}
\end{itemize}


\begin{lstlisting}[language=typescript]
export type AcpSessionMeta = {
  backend: string;
  agent: string;
  runtimeSessionName: string;
  mode: "persistent" | "oneshot";
  cwd?: string;
  state: "idle" | "running" | "error";
  lastActivityAt: number;
  lastError?: string;
};
\end{lstlisting}


Storage rules:

\begin{itemize}
  \item keep \texttt{SessionEntry.acp} as a compatibility projection during migration
  \item process ids and sockets stay in memory only
  \item durable lifecycle and run status live in ACP DB, not generic session JSON
  \item if runtime owner dies, gateway rehydrates from ACP DB and resumes from checkpoints
\end{itemize}


\subsubsection{Routing and delivery}


Inbound:

\begin{itemize}
  \item keep current thread binding lookup as first routing step
  \item if bound target is ACP session, route to ACP runtime branch instead of \texttt{getReplyFromConfig}
  \item explicit \texttt{/acp steer} command uses \texttt{mode: "steer"}
\end{itemize}


Outbound:

\begin{itemize}
  \item ACP event stream is normalized to OpenClaw reply chunks
  \item delivery target is resolved through existing bound destination path
  \item when a bound thread is active for that session turn, parent channel completion is suppressed
\end{itemize}


Streaming policy:

\begin{itemize}
  \item stream partial output with coalescing window
  \item configurable min interval and max chunk bytes to stay under Discord rate limits
  \item final message always emitted on completion or failure
\end{itemize}


\subsubsection{State machines and transaction boundaries}


Session state machine:

\begin{itemize}
  \item \texttt{creating -> idle -> running -> idle}
  \item \texttt{running -> cancelling -> idle | error}
  \item \texttt{idle -> closed}
  \item \texttt{error -> idle | closed}
\end{itemize}


Run state machine:

\begin{itemize}
  \item \texttt{queued -> running -> completed}
  \item \texttt{running -> failed | cancelled}
  \item \texttt{queued -> cancelled}
\end{itemize}


Required transaction boundaries:

\begin{itemize}
  \item spawn transaction
  \item create ACP session row
  \item create/update ACP thread binding row
  \item enqueue initial run row
  \item close transaction
  \item mark session closed
  \item delete/expire binding rows
  \item write final close event
  \item cancel transaction
  \item mark target run cancelling/cancelled with idempotency key
\end{itemize}


No partial success is allowed across these boundaries.

\subsubsection{Per-session actor model}


\texttt{AcpSessionManager} runs one actor per ACP session key:

\begin{itemize}
  \item actor mailbox serializes \texttt{submit}, \texttt{cancel}, \texttt{close}, and \texttt{stream} side effects
  \item actor owns runtime handle hydration and runtime adapter process lifecycle for that session
  \item actor writes run events in-order (\texttt{seq}) before any Discord delivery
  \item actor updates delivery checkpoints after successful outbound send
\end{itemize}


This removes cross-turn races and prevents duplicate or out-of-order thread output.

\subsubsection{Idempotency and delivery projection}


All external ACP actions must carry idempotency keys:

\begin{itemize}
  \item spawn idempotency key
  \item prompt/steer idempotency key
  \item cancel idempotency key
  \item close idempotency key
\end{itemize}


Delivery rules:

\begin{itemize}
  \item Discord messages are derived from \texttt{acp\_events} plus \texttt{acp\_delivery\_checkpoint}
  \item retries resume from checkpoint without re-sending already delivered chunks
  \item final reply emission is exactly-once per run from projection logic
\end{itemize}


\subsubsection{Recovery and self-healing}


On gateway start:

\begin{itemize}
  \item load non-terminal ACP sessions (\texttt{creating}, \texttt{idle}, \texttt{running}, \texttt{cancelling}, \texttt{error})
  \item recreate actors lazily on first inbound event or eagerly under configured cap
  \item reconcile any \texttt{running} runs missing heartbeats and mark \texttt{failed} or recover via adapter
\end{itemize}


On inbound Discord thread message:

\begin{itemize}
  \item if binding exists but ACP session is missing, fail closed with explicit stale-binding message
  \item optionally auto-unbind stale binding after operator-safe validation
  \item never silently route stale ACP bindings to normal LLM path
\end{itemize}


\subsubsection{Lifecycle and safety}


Supported operations:

\begin{itemize}
  \item cancel current run: \texttt{/acp cancel}
  \item unbind thread: \texttt{/unfocus}
  \item close ACP session: \texttt{/acp close}
  \item auto close idle sessions by effective TTL
\end{itemize}


TTL policy:

\begin{itemize}
  \item effective TTL is minimum of
  \item global/session TTL
  \item Discord thread binding TTL
  \item ACP runtime owner TTL
\end{itemize}


Safety controls:

\begin{itemize}
  \item allowlist ACP agents by name
  \item restrict workspace roots for ACP sessions
  \item env allowlist passthrough
  \item max concurrent ACP sessions per account and globally
  \item bounded restart backoff for runtime crashes
\end{itemize}


\subsection{Config surface}


Core keys:

\begin{itemize}
  \item \texttt{acp.enabled}
  \item \texttt{acp.dispatch.enabled} (independent ACP routing kill switch)
  \item \texttt{acp.backend} (default \texttt{acpx})
  \item \texttt{acp.defaultAgent}
  \item \texttt{acp.allowedAgents[]}
  \item \texttt{acp.maxConcurrentSessions}
  \item \texttt{acp.stream.coalesceIdleMs}
  \item \texttt{acp.stream.maxChunkChars}
  \item \texttt{acp.runtime.ttlMinutes}
  \item \texttt{acp.controlPlane.store} (\texttt{sqlite} default)
  \item \texttt{acp.controlPlane.storePath}
  \item \texttt{acp.controlPlane.recovery.eagerActors}
  \item \texttt{acp.controlPlane.recovery.reconcileRunningAfterMs}
  \item \texttt{acp.controlPlane.checkpoint.flushEveryEvents}
  \item \texttt{acp.controlPlane.checkpoint.flushEveryMs}
  \item \texttt{acp.idempotency.ttlHours}
  \item \texttt{channels.discord.threadBindings.spawnAcpSessions}
\end{itemize}


Plugin/backend keys (acpx plugin section):

\begin{itemize}
  \item backend command/path overrides
  \item backend env allowlist
  \item backend per-agent presets
  \item backend startup/stop timeouts
  \item backend max inflight runs per session
\end{itemize}


\subsection{Implementation specification}


\subsubsection{Control-plane modules (new)}


Add dedicated ACP control-plane modules in core:

\begin{itemize}
  \item \texttt{src/acp/control-plane/manager.ts}
  \item owns ACP actors, lifecycle transitions, command serialization
  \item \texttt{src/acp/control-plane/store.ts}
  \item SQLite schema management, transactions, query helpers
  \item \texttt{src/acp/control-plane/events.ts}
  \item typed ACP event definitions and serialization
  \item \texttt{src/acp/control-plane/checkpoint.ts}
  \item durable delivery checkpoints and replay cursors
  \item \texttt{src/acp/control-plane/idempotency.ts}
  \item idempotency key reservation and response replay
  \item \texttt{src/acp/control-plane/recovery.ts}
  \item boot-time reconciliation and actor rehydrate plan
\end{itemize}


Compatibility bridge modules:

\begin{itemize}
  \item \texttt{src/acp/runtime/session-meta.ts}
  \item remains temporarily for projection into \texttt{SessionEntry.acp}
  \item must stop being source-of-truth after migration cutover
\end{itemize}


\subsubsection{Required invariants (must enforce in code)}


\begin{itemize}
  \item ACP session creation and thread bind are atomic (single transaction)
  \item there is at most one active run per ACP session actor at a time
  \item event \texttt{seq} is strictly increasing per run
  \item delivery checkpoint never advances past last committed event
  \item idempotency replay returns previous success payload for duplicate command keys
  \item stale/missing ACP metadata cannot route into normal non-ACP reply path
\end{itemize}


\subsubsection{Core touchpoints}


Core files to change:

\begin{itemize}
  \item \texttt{src/auto-reply/reply/dispatch-from-config.ts}
  \item ACP branch calls \texttt{AcpSessionManager.submit} and event-projection delivery
  \item remove direct ACP fallback that bypasses control-plane invariants
  \item \texttt{src/auto-reply/reply/inbound-context.ts} (or nearest normalized context boundary)
  \item expose normalized routing keys and idempotency seeds for ACP control plane
  \item \texttt{src/config/sessions/types.ts}
  \item keep \texttt{SessionEntry.acp} as projection-only compatibility field
  \item \texttt{src/gateway/server-methods/sessions.ts}
  \item reset/delete/archive must call ACP manager close/unbind transaction path
  \item \texttt{src/infra/outbound/bound-delivery-router.ts}
  \item enforce fail-closed destination behavior for ACP bound session turns
  \item \texttt{src/discord/monitor/thread-bindings.ts}
  \item add ACP stale-binding validation helpers wired to control-plane lookups
  \item \texttt{src/auto-reply/reply/commands-acp.ts}
  \item route spawn/cancel/close/steer through ACP manager APIs
  \item \texttt{src/agents/acp-spawn.ts}
  \item stop ad-hoc metadata writes; call ACP manager spawn transaction
  \item \texttt{src/plugin-sdk/**} and plugin runtime bridge
  \item expose ACP backend registration and health semantics cleanly
\end{itemize}


Core files explicitly not replaced:

\begin{itemize}
  \item \texttt{src/discord/monitor/message-handler.preflight.ts}
  \item keep thread binding override behavior as the canonical session-key resolver
\end{itemize}


\subsubsection{ACP runtime registry API}


Add a core registry module:

\begin{itemize}
  \item \texttt{src/acp/runtime/registry.ts}
\end{itemize}


Required API:

\begin{lstlisting}[language=typescript]
export type AcpRuntimeBackend = {
  id: string;
  runtime: AcpRuntime;
  healthy?: () => boolean;
};

export function registerAcpRuntimeBackend(backend: AcpRuntimeBackend): void;
export function unregisterAcpRuntimeBackend(id: string): void;
export function getAcpRuntimeBackend(id?: string): AcpRuntimeBackend | null;
export function requireAcpRuntimeBackend(id?: string): AcpRuntimeBackend;
\end{lstlisting}


Behavior:

\begin{itemize}
  \item \texttt{requireAcpRuntimeBackend} throws a typed ACP backend missing error when unavailable
  \item plugin service registers backend on \texttt{start} and unregisters on \texttt{stop}
  \item runtime lookups are read-only and process-local
\end{itemize}


\subsubsection{acpx runtime plugin contract (implementation detail)}


For the first production backend (\texttt{extensions/acpx}), OpenClaw and acpx are
connected with a strict command contract:

\begin{itemize}
  \item backend id: \texttt{acpx}
  \item plugin service id: \texttt{acpx-runtime}
  \item runtime handle encoding: \texttt{runtimeSessionName = acpx:v1:}
  \item encoded payload fields:
  \item \texttt{name} (acpx named session; uses OpenClaw \texttt{sessionKey})
  \item \texttt{agent} (acpx agent command)
  \item \texttt{cwd} (session workspace root)
  \item \texttt{mode} (\texttt{persistent | oneshot})
\end{itemize}


Command mapping:

\begin{itemize}
  \item ensure session:
  \item \texttt{acpx --format json --json-strict --cwd   sessions ensure --name }
  \item prompt turn:
  \item \texttt{acpx --format json --json-strict --cwd   prompt --session  --file -}
  \item cancel:
  \item \texttt{acpx --format json --json-strict --cwd   cancel --session }
  \item close:
  \item \texttt{acpx --format json --json-strict --cwd   sessions close }
\end{itemize}


Streaming:

\begin{itemize}
  \item OpenClaw consumes ndjson events from \texttt{acpx --format json --json-strict}
  \item \texttt{text} => \texttt{text\_delta/output}
  \item \texttt{thought} => \texttt{text\_delta/thought}
  \item \texttt{tool\_call} => \texttt{tool\_call}
  \item \texttt{done} => \texttt{done}
  \item \texttt{error} => \texttt{error}
\end{itemize}


\subsubsection{Session schema patch}


Patch \texttt{SessionEntry} in \texttt{src/config/sessions/types.ts}:

\begin{lstlisting}[language=typescript]
type SessionAcpMeta = {
  backend: string;
  agent: string;
  runtimeSessionName: string;
  mode: "persistent" | "oneshot";
  cwd?: string;
  state: "idle" | "running" | "error";
  lastActivityAt: number;
  lastError?: string;
};
\end{lstlisting}


Persisted field:

\begin{itemize}
  \item \texttt{SessionEntry.acp?: SessionAcpMeta}
\end{itemize}


Migration rules:

\begin{itemize}
  \item phase A: dual-write (\texttt{acp} projection + ACP SQLite source-of-truth)
  \item phase B: read-primary from ACP SQLite, fallback-read from legacy \texttt{SessionEntry.acp}
  \item phase C: migration command backfills missing ACP rows from valid legacy entries
  \item phase D: remove fallback-read and keep projection optional for UX only
  \item legacy fields (\texttt{cliSessionIds}, \texttt{claudeCliSessionId}) remain untouched
\end{itemize}


\subsubsection{Error contract}


Add stable ACP error codes and user-facing messages:

\begin{itemize}
  \item \texttt{ACP\_BACKEND\_MISSING}
  \item message: \texttt{ACP runtime backend is not configured. Install and enable the acpx runtime plugin.}
  \item \texttt{ACP\_BACKEND\_UNAVAILABLE}
  \item message: \texttt{ACP runtime backend is currently unavailable. Try again in a moment.}
  \item \texttt{ACP\_SESSION\_INIT\_FAILED}
  \item message: \texttt{Could not initialize ACP session runtime.}
  \item \texttt{ACP\_TURN\_FAILED}
  \item message: \texttt{ACP turn failed before completion.}
\end{itemize}


Rules:

\begin{itemize}
  \item return actionable user-safe message in-thread
  \item log detailed backend/system error only in runtime logs
  \item never silently fall back to normal LLM path when ACP routing was explicitly selected
\end{itemize}


\subsubsection{Duplicate delivery arbitration}


Single routing rule for ACP bound turns:

\begin{itemize}
  \item if an active thread binding exists for the target ACP session and requester context, deliver only to that bound thread
  \item do not also send to parent channel for the same turn
  \item if bound destination selection is ambiguous, fail closed with explicit error (no implicit parent fallback)
  \item if no active binding exists, use normal session destination behavior
\end{itemize}


\subsubsection{Observability and operational readiness}


Required metrics:

\begin{itemize}
  \item ACP spawn success/failure count by backend and error code
  \item ACP run latency percentiles (queue wait, runtime turn time, delivery projection time)
  \item ACP actor restart count and restart reason
  \item stale-binding detection count
  \item idempotency replay hit rate
  \item Discord delivery retry and rate-limit counters
\end{itemize}


Required logs:

\begin{itemize}
  \item structured logs keyed by \texttt{sessionKey}, \texttt{runId}, \texttt{backend}, \texttt{threadId}, \texttt{idempotencyKey}
  \item explicit state transition logs for session and run state machines
  \item adapter command logs with redaction-safe arguments and exit summary
\end{itemize}


Required diagnostics:

\begin{itemize}
  \item \texttt{/acp sessions} includes state, active run, last error, and binding status
  \item \texttt{/acp doctor} (or equivalent) validates backend registration, store health, and stale bindings
\end{itemize}


\subsubsection{Config precedence and effective values}


ACP enablement precedence:

\begin{itemize}
  \item account override: \texttt{channels.discord.accounts..threadBindings.spawnAcpSessions}
  \item channel override: \texttt{channels.discord.threadBindings.spawnAcpSessions}
  \item global ACP gate: \texttt{acp.enabled}
  \item dispatch gate: \texttt{acp.dispatch.enabled}
  \item backend availability: registered backend for \texttt{acp.backend}
\end{itemize}


Auto-enable behavior:

\begin{itemize}
  \item when ACP is configured (\texttt{acp.enabled=true}, \texttt{acp.dispatch.enabled=true}, or
\end{itemize}

\texttt{acp.backend=acpx}), plugin auto-enable marks \texttt{plugins.entries.acpx.enabled=true}
unless denylisted or explicitly disabled

TTL effective value:

\begin{itemize}
  \item \texttt{min(session ttl, discord thread binding ttl, acp runtime ttl)}
\end{itemize}


\subsubsection{Test map}


Unit tests:

\begin{itemize}
  \item \texttt{src/acp/runtime/registry.test.ts} (new)
  \item \texttt{src/auto-reply/reply/dispatch-from-config.acp.test.ts} (new)
  \item \texttt{src/infra/outbound/bound-delivery-router.test.ts} (extend ACP fail-closed cases)
  \item \texttt{src/config/sessions/types.test.ts} or nearest session-store tests (ACP metadata persistence)
\end{itemize}


Integration tests:

\begin{itemize}
  \item \texttt{src/discord/monitor/reply-delivery.test.ts} (bound ACP delivery target behavior)
  \item \texttt{src/discord/monitor/message-handler.preflight*.test.ts} (bound ACP session-key routing continuity)
  \item acpx plugin runtime tests in backend package (service register/start/stop + event normalization)
\end{itemize}


Gateway e2e tests:

\begin{itemize}
  \item \texttt{src/gateway/server.sessions.gateway-server-sessions-a.e2e.test.ts} (extend ACP reset/delete lifecycle coverage)
  \item ACP thread turn roundtrip e2e for spawn, message, stream, cancel, unfocus, restart recovery
\end{itemize}


\subsubsection{Rollout guard}


Add independent ACP dispatch kill switch:

\begin{itemize}
  \item \texttt{acp.dispatch.enabled} default \texttt{false} for first release
  \item when disabled:
  \item ACP spawn/focus control commands may still bind sessions
  \item ACP dispatch path does not activate
  \item user receives explicit message that ACP dispatch is disabled by policy
  \item after canary validation, default can be flipped to \texttt{true} in a later release
\end{itemize}


\subsection{Command and UX plan}


\subsubsection{New commands}


\begin{itemize}
  \item \texttt{/acp spawn  [--mode persistent|oneshot] [--thread auto|here|off]}
  \item \texttt{/acp cancel [session]}
  \item \texttt{/acp steer }
  \item \texttt{/acp close [session]}
  \item \texttt{/acp sessions}
\end{itemize}


\subsubsection{Existing command compatibility}


\begin{itemize}
  \item \texttt{/focus } continues to support ACP targets
  \item \texttt{/unfocus} keeps current semantics
  \item \texttt{/session idle} and \texttt{/session max-age} replace the old TTL override
\end{itemize}


\subsection{Phased rollout}


\subsubsection{Phase 0 ADR and schema freeze}


\begin{itemize}
  \item ship ADR for ACP control-plane ownership and adapter boundaries
  \item freeze DB schema (\texttt{acp\_sessions}, \texttt{acp\_runs}, \texttt{acp\_bindings}, \texttt{acp\_events}, \texttt{acp\_delivery\_checkpoint}, \texttt{acp\_idempotency})
  \item define stable ACP error codes, event contract, and state-transition guards
\end{itemize}


\subsubsection{Phase 1 Control-plane foundation in core}


\begin{itemize}
  \item implement \texttt{AcpSessionManager} and per-session actor runtime
  \item implement ACP SQLite store and transaction helpers
  \item implement idempotency store and replay helpers
  \item implement event append + delivery checkpoint modules
  \item wire spawn/cancel/close APIs to manager with transactional guarantees
\end{itemize}


\subsubsection{Phase 2 Core routing and lifecycle integration}


\begin{itemize}
  \item route thread-bound ACP turns from dispatch pipeline into ACP manager
  \item enforce fail-closed routing when ACP binding/session invariants fail
  \item integrate reset/delete/archive/unfocus lifecycle with ACP close/unbind transactions
  \item add stale-binding detection and optional auto-unbind policy
\end{itemize}


\subsubsection{Phase 3 acpx backend adapter/plugin}


\begin{itemize}
  \item implement \texttt{acpx} adapter against runtime contract (\texttt{ensureSession}, \texttt{submit}, \texttt{stream}, \texttt{cancel}, \texttt{close})
  \item add backend health checks and startup/teardown registration
  \item normalize acpx ndjson events into ACP runtime events
  \item enforce backend timeouts, process supervision, and restart/backoff policy
\end{itemize}


\subsubsection{Phase 4 Delivery projection and channel UX (Discord first)}


\begin{itemize}
  \item implement event-driven channel projection with checkpoint resume (Discord first)
  \item coalesce streaming chunks with rate-limit aware flush policy
  \item guarantee exactly-once final completion message per run
  \item ship \texttt{/acp spawn}, \texttt{/acp cancel}, \texttt{/acp steer}, \texttt{/acp close}, \texttt{/acp sessions}
\end{itemize}


\subsubsection{Phase 5 Migration and cutover}


\begin{itemize}
  \item introduce dual-write to \texttt{SessionEntry.acp} projection plus ACP SQLite source-of-truth
  \item add migration utility for legacy ACP metadata rows
  \item flip read path to ACP SQLite primary
  \item remove legacy fallback routing that depends on missing \texttt{SessionEntry.acp}
\end{itemize}


\subsubsection{Phase 6 Hardening, SLOs, and scale limits}


\begin{itemize}
  \item enforce concurrency limits (global/account/session), queue policies, and timeout budgets
  \item add full telemetry, dashboards, and alert thresholds
  \item chaos-test crash recovery and duplicate-delivery suppression
  \item publish runbook for backend outage, DB corruption, and stale-binding remediation
\end{itemize}


\subsubsection{Full implementation checklist}


\begin{itemize}
  \item core control-plane modules and tests
  \item DB migrations and rollback plan
  \item ACP manager API integration across dispatch and commands
  \item adapter registration interface in plugin runtime bridge
  \item acpx adapter implementation and tests
  \item thread-capable channel delivery projection logic with checkpoint replay (Discord first)
  \item lifecycle hooks for reset/delete/archive/unfocus
  \item stale-binding detector and operator-facing diagnostics
  \item config validation and precedence tests for all new ACP keys
  \item operational docs and troubleshooting runbook
\end{itemize}


\subsection{Test plan}


Unit tests:

\begin{itemize}
  \item ACP DB transaction boundaries (spawn/bind/enqueue atomicity, cancel, close)
  \item ACP state-machine transition guards for sessions and runs
  \item idempotency reservation/replay semantics across all ACP commands
  \item per-session actor serialization and queue ordering
  \item acpx event parser and chunk coalescer
  \item runtime supervisor restart and backoff policy
  \item config precedence and effective TTL calculation
  \item core ACP routing branch selection and fail-closed behavior when backend/session is invalid
\end{itemize}


Integration tests:

\begin{itemize}
  \item fake ACP adapter process for deterministic streaming and cancel behavior
  \item ACP manager + dispatch integration with transactional persistence
  \item thread-bound inbound routing to ACP session key
  \item thread-bound outbound delivery suppresses parent channel duplication
  \item checkpoint replay recovers after delivery failure and resumes from last event
  \item plugin service registration and teardown of ACP runtime backend
\end{itemize}


Gateway e2e tests:

\begin{itemize}
  \item spawn ACP with thread, exchange multi-turn prompts, unfocus
  \item gateway restart with persisted ACP DB and bindings, then continue same session
  \item concurrent ACP sessions in multiple threads have no cross-talk
  \item duplicate command retries (same idempotency key) do not create duplicate runs or replies
  \item stale-binding scenario yields explicit error and optional auto-clean behavior
\end{itemize}


\subsection{Risks and mitigations}


\begin{itemize}
  \item Duplicate deliveries during transition
  \item Mitigation: single destination resolver and idempotent event checkpoint
  \item Runtime process churn under load
  \item Mitigation: long lived per session owners + concurrency caps + backoff
  \item Plugin absent or misconfigured
  \item Mitigation: explicit operator-facing error and fail-closed ACP routing (no implicit fallback to normal session path)
  \item Config confusion between subagent and ACP gates
  \item Mitigation: explicit ACP keys and command feedback that includes effective policy source
  \item Control-plane store corruption or migration bugs
  \item Mitigation: WAL mode, backup/restore hooks, migration smoke tests, and read-only fallback diagnostics
  \item Actor deadlocks or mailbox starvation
  \item Mitigation: watchdog timers, actor health probes, and bounded mailbox depth with rejection telemetry
\end{itemize}


\subsection{Acceptance checklist}


\begin{itemize}
  \item ACP session spawn can create or bind a thread in a supported channel adapter (currently Discord)
  \item all thread messages route to bound ACP session only
  \item ACP outputs appear in the same thread identity with streaming or batches
  \item no duplicate output in parent channel for bound turns
  \item spawn+bind+initial enqueue are atomic in persistent store
  \item ACP command retries are idempotent and do not duplicate runs or outputs
  \item cancel, close, unfocus, archive, reset, and delete perform deterministic cleanup
  \item crash restart preserves mapping and resumes multi turn continuity
  \item concurrent thread bound ACP sessions work independently
  \item ACP backend missing state produces clear actionable error
  \item stale bindings are detected and surfaced explicitly (with optional safe auto-clean)
  \item control-plane metrics and diagnostics are available for operators
  \item new unit, integration, and e2e coverage passes
\end{itemize}


\subsection{Addendum: targeted refactors for current implementation (status)}


These are non-blocking follow-ups to keep the ACP path maintainable after the current feature set lands.

\subsubsection{1) Centralize ACP dispatch policy evaluation (completed)}


\begin{itemize}
  \item implemented via shared ACP policy helpers in \texttt{src/acp/policy.ts}
  \item dispatch, ACP command lifecycle handlers, and ACP spawn path now consume shared policy logic
\end{itemize}


\subsubsection{2) Split ACP command handler by subcommand domain (completed)}


\begin{itemize}
  \item \texttt{src/auto-reply/reply/commands-acp.ts} is now a thin router
  \item subcommand behavior is split into:
  \item \texttt{src/auto-reply/reply/commands-acp/lifecycle.ts}
  \item \texttt{src/auto-reply/reply/commands-acp/runtime-options.ts}
  \item \texttt{src/auto-reply/reply/commands-acp/diagnostics.ts}
  \item shared helpers in \texttt{src/auto-reply/reply/commands-acp/shared.ts}
\end{itemize}


\subsubsection{3) Split ACP session manager by responsibility (completed)}


\begin{itemize}
  \item manager is split into:
  \item \texttt{src/acp/control-plane/manager.ts} (public facade + singleton)
  \item \texttt{src/acp/control-plane/manager.core.ts} (manager implementation)
  \item \texttt{src/acp/control-plane/manager.types.ts} (manager types/deps)
  \item \texttt{src/acp/control-plane/manager.utils.ts} (normalization + helper functions)
\end{itemize}


\subsubsection{4) Optional acpx runtime adapter cleanup}


\begin{itemize}
  \item \texttt{extensions/acpx/src/runtime.ts} can be split into:
  \item process execution/supervision
  \item ndjson event parsing/normalization
  \item runtime API surface (\texttt{submit}, \texttt{cancel}, \texttt{close}, etc.)
  \item improves testability and makes backend behavior easier to audit
\end{itemize}



\clearpage

% === 源文件: experiments/plans/acp-unified-streaming-refactor.md ===
\section{Unified Runtime Streaming Refactor Plan}


\subsection{Objective}


Deliver one shared streaming pipeline for \texttt{main}, \texttt{subagent}, and \texttt{acp} so all runtimes get identical coalescing, chunking, delivery ordering, and crash recovery behavior.

\subsection{Why this exists}


\begin{itemize}
  \item Current behavior is split across multiple runtime-specific shaping paths.
  \item Formatting/coalescing bugs can be fixed in one path but remain in others.
  \item Delivery consistency, duplicate suppression, and recovery semantics are harder to reason about.
\end{itemize}


\subsection{Target architecture}


Single pipeline, runtime-specific adapters:

\begin{enumerate}
  \item Runtime adapters emit canonical events only.
  \item Shared stream assembler coalesces and finalizes text/tool/status events.
  \item Shared channel projector applies channel-specific chunking/formatting once.
  \item Shared delivery ledger enforces idempotent send/replay semantics.
  \item Outbound channel adapter executes sends and records delivery checkpoints.
\end{enumerate}


Canonical event contract:

\begin{itemize}
  \item \texttt{turn\_started}
  \item \texttt{text\_delta}
  \item \texttt{block\_final}
  \item \texttt{tool\_started}
  \item \texttt{tool\_finished}
  \item \texttt{status}
  \item \texttt{turn\_completed}
  \item \texttt{turn\_failed}
  \item \texttt{turn\_cancelled}
\end{itemize}


\subsection{Workstreams}


\subsubsection{1) Canonical streaming contract}


\begin{itemize}
  \item Define strict event schema + validation in core.
  \item Add adapter contract tests to guarantee each runtime emits compatible events.
  \item Reject malformed runtime events early and surface structured diagnostics.
\end{itemize}


\subsubsection{2) Shared stream processor}


\begin{itemize}
  \item Replace runtime-specific coalescer/projector logic with one processor.
  \item Processor owns text delta buffering, idle flush, max-chunk splitting, and completion flush.
  \item Move ACP/main/subagent config resolution into one helper to prevent drift.
\end{itemize}


\subsubsection{3) Shared channel projection}


\begin{itemize}
  \item Keep channel adapters dumb: accept finalized blocks and send.
  \item Move Discord-specific chunking quirks to channel projector only.
  \item Keep pipeline channel-agnostic before projection.
\end{itemize}


\subsubsection{4) Delivery ledger + replay}


\begin{itemize}
  \item Add per-turn/per-chunk delivery IDs.
  \item Record checkpoints before and after physical send.
  \item On restart, replay pending chunks idempotently and avoid duplicates.
\end{itemize}


\subsubsection{5) Migration and cutover}


\begin{itemize}
  \item Phase 1: shadow mode (new pipeline computes output but old path sends; compare).
  \item Phase 2: runtime-by-runtime cutover (\texttt{acp}, then \texttt{subagent}, then \texttt{main} or reverse by risk).
  \item Phase 3: delete legacy runtime-specific streaming code.
\end{itemize}


\subsection{Non-goals}


\begin{itemize}
  \item No changes to ACP policy/permissions model in this refactor.
  \item No channel-specific feature expansion outside projection compatibility fixes.
  \item No transport/backend redesign (acpx plugin contract remains as-is unless needed for event parity).
\end{itemize}


\subsection{Risks and mitigations}


\begin{itemize}
  \item Risk: behavioral regressions in existing main/subagent paths.
\end{itemize}

Mitigation: shadow mode diffing + adapter contract tests + channel e2e tests.
\begin{itemize}
  \item Risk: duplicate sends during crash recovery.
\end{itemize}

Mitigation: durable delivery IDs + idempotent replay in delivery adapter.
\begin{itemize}
  \item Risk: runtime adapters diverge again.
\end{itemize}

Mitigation: required shared contract test suite for all adapters.

\subsection{Acceptance criteria}


\begin{itemize}
  \item All runtimes pass shared streaming contract tests.
  \item Discord ACP/main/subagent produce equivalent spacing/chunking behavior for tiny deltas.
  \item Crash/restart replay sends no duplicate chunk for the same delivery ID.
  \item Legacy ACP projector/coalescer path is removed.
  \item Streaming config resolution is shared and runtime-independent.
\end{itemize}



\clearpage

% === 源文件: experiments/plans/browser-evaluate-cdp-refactor.md ===
\section{Browser Evaluate CDP Refactor Plan}


\subsection{Context}


\texttt{act:evaluate} executes user provided JavaScript in the page. Today it runs via Playwright
(\texttt{page.evaluate} or \texttt{locator.evaluate}). Playwright serializes CDP commands per page, so a
stuck or long running evaluate can block the page command queue and make every later action
on that tab look "stuck".

PR \#13498 adds a pragmatic safety net (bounded evaluate, abort propagation, and best-effort
recovery). This document describes a larger refactor that makes \texttt{act:evaluate} inherently
isolated from Playwright so a stuck evaluate cannot wedge normal Playwright operations.

\subsection{Goals}


\begin{itemize}
  \item \texttt{act:evaluate} cannot permanently block later browser actions on the same tab.
  \item Timeouts are single source of truth end to end so a caller can rely on a budget.
  \item Abort and timeout are treated the same way across HTTP and in-process dispatch.
  \item Element targeting for evaluate is supported without switching everything off Playwright.
  \item Maintain backward compatibility for existing callers and payloads.
\end{itemize}


\subsection{Non-goals}


\begin{itemize}
  \item Replace all browser actions (click, type, wait, etc.) with CDP implementations.
  \item Remove the existing safety net introduced in PR \#13498 (it remains a useful fallback).
  \item Introduce new unsafe capabilities beyond the existing \texttt{browser.evaluateEnabled} gate.
  \item Add process isolation (worker process/thread) for evaluate. If we still see hard to recover
\end{itemize}

stuck states after this refactor, that is a follow-up idea.

\subsection{Current Architecture (Why It Gets Stuck)}


At a high level:

\begin{itemize}
  \item Callers send \texttt{act:evaluate} to the browser control service.
  \item The route handler calls into Playwright to execute the JavaScript.
  \item Playwright serializes page commands, so an evaluate that never finishes blocks the queue.
  \item A stuck queue means later click/type/wait operations on the tab can appear to hang.
\end{itemize}


\subsection{Proposed Architecture}


\subsubsection{1. Deadline Propagation}


Introduce a single budget concept and derive everything from it:

\begin{itemize}
  \item Caller sets \texttt{timeoutMs} (or a deadline in the future).
  \item The outer request timeout, route handler logic, and the execution budget inside the page
\end{itemize}

all use the same budget, with small headroom where needed for serialization overhead.
\begin{itemize}
  \item Abort is propagated as an \texttt{AbortSignal} everywhere so cancellation is consistent.
\end{itemize}


Implementation direction:

\begin{itemize}
  \item Add a small helper (for example \texttt{createBudget(\{ timeoutMs, signal \})}) that returns:
  \item \texttt{signal}: the linked AbortSignal
  \item \texttt{deadlineAtMs}: absolute deadline
  \item \texttt{remainingMs()}: remaining budget for child operations
  \item Use this helper in:
  \item \texttt{src/browser/client-fetch.ts} (HTTP and in-process dispatch)
  \item \texttt{src/node-host/runner.ts} (proxy path)
  \item browser action implementations (Playwright and CDP)
\end{itemize}


\subsubsection{2. Separate Evaluate Engine (CDP Path)}


Add a CDP based evaluate implementation that does not share Playwright's per page command
queue. The key property is that the evaluate transport is a separate WebSocket connection
and a separate CDP session attached to the target.

Implementation direction:

\begin{itemize}
  \item New module, for example \texttt{src/browser/cdp-evaluate.ts}, that:
  \item Connects to the configured CDP endpoint (browser level socket).
  \item Uses \texttt{Target.attachToTarget(\{ targetId, flatten: true \})} to get a \texttt{sessionId}.
  \item Runs either:
  \item \texttt{Runtime.evaluate} for page level evaluate, or
  \item \texttt{DOM.resolveNode} plus \texttt{Runtime.callFunctionOn} for element evaluate.
  \item On timeout or abort:
  \item Sends \texttt{Runtime.terminateExecution} best-effort for the session.
  \item Closes the WebSocket and returns a clear error.
\end{itemize}


Notes:

\begin{itemize}
  \item This still executes JavaScript in the page, so termination can have side effects. The win
\end{itemize}

is that it does not wedge the Playwright queue, and it is cancelable at the transport
layer by killing the CDP session.

\subsubsection{3. Ref Story (Element Targeting Without A Full Rewrite)}


The hard part is element targeting. CDP needs a DOM handle or \texttt{backendDOMNodeId}, while
today most browser actions use Playwright locators based on refs from snapshots.

Recommended approach: keep existing refs, but attach an optional CDP resolvable id.

\paragraph{3.1 Extend Stored Ref Info}


Extend the stored role ref metadata to optionally include a CDP id:

\begin{itemize}
  \item Today: \texttt{\{ role, name, nth \}}
  \item Proposed: \texttt{\{ role, name, nth, backendDOMNodeId?: number \}}
\end{itemize}


This keeps all existing Playwright based actions working and allows CDP evaluate to accept
the same \texttt{ref} value when the \texttt{backendDOMNodeId} is available.

\paragraph{3.2 Populate backendDOMNodeId At Snapshot Time}


When producing a role snapshot:

\begin{enumerate}
  \item Generate the existing role ref map as today (role, name, nth).
  \item Fetch the AX tree via CDP (\texttt{Accessibility.getFullAXTree}) and compute a parallel map of
\end{enumerate}

\texttt{(role, name, nth) -> backendDOMNodeId} using the same duplicate handling rules.
\begin{enumerate}
  \item Merge the id back into the stored ref info for the current tab.
\end{enumerate}


If mapping fails for a ref, leave \texttt{backendDOMNodeId} undefined. This makes the feature
best-effort and safe to roll out.

\paragraph{3.3 Evaluate Behavior With Ref}


In \texttt{act:evaluate}:

\begin{itemize}
  \item If \texttt{ref} is present and has \texttt{backendDOMNodeId}, run element evaluate via CDP.
  \item If \texttt{ref} is present but has no \texttt{backendDOMNodeId}, fall back to the Playwright path (with
\end{itemize}

the safety net).

Optional escape hatch:

\begin{itemize}
  \item Extend the request shape to accept \texttt{backendDOMNodeId} directly for advanced callers (and
\end{itemize}

for debugging), while keeping \texttt{ref} as the primary interface.

\subsubsection{4. Keep A Last Resort Recovery Path}


Even with CDP evaluate, there are other ways to wedge a tab or a connection. Keep the
existing recovery mechanisms (terminate execution + disconnect Playwright) as a last resort
for:

\begin{itemize}
  \item legacy callers
  \item environments where CDP attach is blocked
  \item unexpected Playwright edge cases
\end{itemize}


\subsection{Implementation Plan (Single Iteration)}


\subsubsection{Deliverables}


\begin{itemize}
  \item A CDP based evaluate engine that runs outside the Playwright per-page command queue.
  \item A single end-to-end timeout/abort budget used consistently by callers and handlers.
  \item Ref metadata that can optionally carry \texttt{backendDOMNodeId} for element evaluate.
  \item \texttt{act:evaluate} prefers the CDP engine when possible and falls back to Playwright when not.
  \item Tests that prove a stuck evaluate does not wedge later actions.
  \item Logs/metrics that make failures and fallbacks visible.
\end{itemize}


\subsubsection{Implementation Checklist}


\begin{enumerate}
  \item Add a shared "budget" helper to link \texttt{timeoutMs} + upstream \texttt{AbortSignal} into:
\end{enumerate}

\begin{itemize}
  \item a single \texttt{AbortSignal}
  \item an absolute deadline
  \item a \texttt{remainingMs()} helper for downstream operations
\end{itemize}

\begin{enumerate}
  \item Update all caller paths to use that helper so \texttt{timeoutMs} means the same thing everywhere:
\end{enumerate}

\begin{itemize}
  \item \texttt{src/browser/client-fetch.ts} (HTTP and in-process dispatch)
  \item \texttt{src/node-host/runner.ts} (node proxy path)
  \item CLI wrappers that call \texttt{/act} (add \texttt{--timeout-ms} to \texttt{browser evaluate})
\end{itemize}

\begin{enumerate}
  \item Implement \texttt{src/browser/cdp-evaluate.ts}:
\end{enumerate}

\begin{itemize}
  \item connect to the browser-level CDP socket
  \item \texttt{Target.attachToTarget} to get a \texttt{sessionId}
  \item run \texttt{Runtime.evaluate} for page evaluate
  \item run \texttt{DOM.resolveNode} + \texttt{Runtime.callFunctionOn} for element evaluate
  \item on timeout/abort: best-effort \texttt{Runtime.terminateExecution} then close the socket
\end{itemize}

\begin{enumerate}
  \item Extend stored role ref metadata to optionally include \texttt{backendDOMNodeId}:
\end{enumerate}

\begin{itemize}
  \item keep existing \texttt{\{ role, name, nth \}} behavior for Playwright actions
  \item add \texttt{backendDOMNodeId?: number} for CDP element targeting
\end{itemize}

\begin{enumerate}
  \item Populate \texttt{backendDOMNodeId} during snapshot creation (best-effort):
\end{enumerate}

\begin{itemize}
  \item fetch AX tree via CDP (\texttt{Accessibility.getFullAXTree})
  \item compute \texttt{(role, name, nth) -> backendDOMNodeId} and merge into the stored ref map
  \item if mapping is ambiguous or missing, leave the id undefined
\end{itemize}

\begin{enumerate}
  \item Update \texttt{act:evaluate} routing:
\end{enumerate}

\begin{itemize}
  \item if no \texttt{ref}: always use CDP evaluate
  \item if \texttt{ref} resolves to a \texttt{backendDOMNodeId}: use CDP element evaluate
  \item otherwise: fall back to Playwright evaluate (still bounded and abortable)
\end{itemize}

\begin{enumerate}
  \item Keep the existing "last resort" recovery path as a fallback, not the default path.
  \item Add tests:
\end{enumerate}

\begin{itemize}
  \item stuck evaluate times out within budget and the next click/type succeeds
  \item abort cancels evaluate (client disconnect or timeout) and unblocks subsequent actions
  \item mapping failures cleanly fall back to Playwright
\end{itemize}

\begin{enumerate}
  \item Add observability:
\end{enumerate}

\begin{itemize}
  \item evaluate duration and timeout counters
  \item terminateExecution usage
  \item fallback rate (CDP -> Playwright) and reasons
\end{itemize}


\subsubsection{Acceptance Criteria}


\begin{itemize}
  \item A deliberately hung \texttt{act:evaluate} returns within the caller budget and does not wedge the
\end{itemize}

tab for later actions.
\begin{itemize}
  \item \texttt{timeoutMs} behaves consistently across CLI, agent tool, node proxy, and in-process calls.
  \item If \texttt{ref} can be mapped to \texttt{backendDOMNodeId}, element evaluate uses CDP; otherwise the
\end{itemize}

fallback path is still bounded and recoverable.

\subsection{Testing Plan}


\begin{itemize}
  \item Unit tests:
  \item \texttt{(role, name, nth)} matching logic between role refs and AX tree nodes.
  \item Budget helper behavior (headroom, remaining time math).
  \item Integration tests:
  \item CDP evaluate timeout returns within budget and does not block the next action.
  \item Abort cancels evaluate and triggers termination best-effort.
  \item Contract tests:
  \item Ensure \texttt{BrowserActRequest} and \texttt{BrowserActResponse} remain compatible.
\end{itemize}


\subsection{Risks And Mitigations}


\begin{itemize}
  \item Mapping is imperfect:
  \item Mitigation: best-effort mapping, fallback to Playwright evaluate, and add debug tooling.
  \item \texttt{Runtime.terminateExecution} has side effects:
  \item Mitigation: only use on timeout/abort and document the behavior in errors.
  \item Extra overhead:
  \item Mitigation: only fetch AX tree when snapshots are requested, cache per target, and keep
\end{itemize}

CDP session short lived.
\begin{itemize}
  \item Extension relay limitations:
  \item Mitigation: use browser level attach APIs when per page sockets are not available, and
\end{itemize}

keep the current Playwright path as fallback.

\subsection{Open Questions}


\begin{itemize}
  \item Should the new engine be configurable as \texttt{playwright}, \texttt{cdp}, or \texttt{auto}?
  \item Do we want to expose a new "nodeRef" format for advanced users, or keep \texttt{ref} only?
  \item How should frame snapshots and selector scoped snapshots participate in AX mapping?
\end{itemize}



\clearpage

% === 源文件: experiments/plans/pty-process-supervision.md ===
\section{PTY and Process Supervision Plan}


\subsection{1. Problem and goal}


We need one reliable lifecycle for long-running command execution across:

\begin{itemize}
  \item \texttt{exec} foreground runs
  \item \texttt{exec} background runs
  \item \texttt{process} follow up actions (\texttt{poll}, \texttt{log}, \texttt{send-keys}, \texttt{paste}, \texttt{submit}, \texttt{kill}, \texttt{remove})
  \item CLI agent runner subprocesses
\end{itemize}


The goal is not just to support PTY. The goal is predictable ownership, cancellation, timeout, and cleanup with no unsafe process matching heuristics.

\subsection{2. Scope and boundaries}


\begin{itemize}
  \item Keep implementation internal in \texttt{src/process/supervisor}.
  \item Do not create a new package for this.
  \item Keep current behavior compatibility where practical.
  \item Do not broaden scope to terminal replay or tmux style session persistence.
\end{itemize}


\subsection{3. Implemented in this branch}


\subsubsection{Supervisor baseline already present}


\begin{itemize}
  \item Supervisor module is in place under \texttt{src/process/supervisor/*}.
  \item Exec runtime and CLI runner are already routed through supervisor spawn and wait.
  \item Registry finalization is idempotent.
\end{itemize}


\subsubsection{This pass completed}


\begin{enumerate}
  \item Explicit PTY command contract
\end{enumerate}


\begin{itemize}
  \item \texttt{SpawnInput} is now a discriminated union in \texttt{src/process/supervisor/types.ts}.
  \item PTY runs require \texttt{ptyCommand} instead of reusing generic \texttt{argv}.
  \item Supervisor no longer rebuilds PTY command strings from argv joins in \texttt{src/process/supervisor/supervisor.ts}.
  \item Exec runtime now passes \texttt{ptyCommand} directly in \texttt{src/agents/bash-tools.exec-runtime.ts}.
\end{itemize}


\begin{enumerate}
  \item Process layer type decoupling
\end{enumerate}


\begin{itemize}
  \item Supervisor types no longer import \texttt{SessionStdin} from agents.
  \item Process local stdin contract lives in \texttt{src/process/supervisor/types.ts} (\texttt{ManagedRunStdin}).
  \item Adapters now depend only on process level types:
  \item \texttt{src/process/supervisor/adapters/child.ts}
  \item \texttt{src/process/supervisor/adapters/pty.ts}
\end{itemize}


\begin{enumerate}
  \item Process tool lifecycle ownership improvement
\end{enumerate}


\begin{itemize}
  \item \texttt{src/agents/bash-tools.process.ts} now requests cancellation through supervisor first.
  \item \texttt{process kill/remove} now use process-tree fallback termination when supervisor lookup misses.
  \item \texttt{remove} keeps deterministic remove behavior by dropping running session entries immediately after termination is requested.
\end{itemize}


\begin{enumerate}
  \item Single source watchdog defaults
\end{enumerate}


\begin{itemize}
  \item Added shared defaults in \texttt{src/agents/cli-watchdog-defaults.ts}.
  \item \texttt{src/agents/cli-backends.ts} consumes the shared defaults.
  \item \texttt{src/agents/cli-runner/reliability.ts} consumes the same shared defaults.
\end{itemize}


\begin{enumerate}
  \item Dead helper cleanup
\end{enumerate}


\begin{itemize}
  \item Removed unused \texttt{killSession} helper path from \texttt{src/agents/bash-tools.shared.ts}.
\end{itemize}


\begin{enumerate}
  \item Direct supervisor path tests added
\end{enumerate}


\begin{itemize}
  \item Added \texttt{src/agents/bash-tools.process.supervisor.test.ts} to cover kill and remove routing through supervisor cancellation.
\end{itemize}


\begin{enumerate}
  \item Reliability gap fixes completed
\end{enumerate}


\begin{itemize}
  \item \texttt{src/agents/bash-tools.process.ts} now falls back to real OS-level process termination when supervisor lookup misses.
  \item \texttt{src/process/supervisor/adapters/child.ts} now uses process-tree termination semantics for default cancel/timeout kill paths.
  \item Added shared process-tree utility in \texttt{src/process/kill-tree.ts}.
\end{itemize}


\begin{enumerate}
  \item PTY contract edge-case coverage added
\end{enumerate}


\begin{itemize}
  \item Added \texttt{src/process/supervisor/supervisor.pty-command.test.ts} for verbatim PTY command forwarding and empty-command rejection.
  \item Added \texttt{src/process/supervisor/adapters/child.test.ts} for process-tree kill behavior in child adapter cancellation.
\end{itemize}


\subsection{4. Remaining gaps and decisions}


\subsubsection{Reliability status}


The two required reliability gaps for this pass are now closed:

\begin{itemize}
  \item \texttt{process kill/remove} now has a real OS termination fallback when supervisor lookup misses.
  \item child cancel/timeout now uses process-tree kill semantics for default kill path.
  \item Regression tests were added for both behaviors.
\end{itemize}


\subsubsection{Durability and startup reconciliation}


Restart behavior is now explicitly defined as in-memory lifecycle only.

\begin{itemize}
  \item \texttt{reconcileOrphans()} remains a no-op in \texttt{src/process/supervisor/supervisor.ts} by design.
  \item Active runs are not recovered after process restart.
  \item This boundary is intentional for this implementation pass to avoid partial persistence risks.
\end{itemize}


\subsubsection{Maintainability follow-ups}


\begin{enumerate}
  \item \texttt{runExecProcess} in \texttt{src/agents/bash-tools.exec-runtime.ts} still handles multiple responsibilities and can be split into focused helpers in a follow-up.
\end{enumerate}


\subsection{5. Implementation plan}


The implementation pass for required reliability and contract items is complete.

Completed:

\begin{itemize}
  \item \texttt{process kill/remove} fallback real termination
  \item process-tree cancellation for child adapter default kill path
  \item regression tests for fallback kill and child adapter kill path
  \item PTY command edge-case tests under explicit \texttt{ptyCommand}
  \item explicit in-memory restart boundary with \texttt{reconcileOrphans()} no-op by design
\end{itemize}


Optional follow-up:

\begin{itemize}
  \item split \texttt{runExecProcess} into focused helpers with no behavior drift
\end{itemize}


\subsection{6. File map}


\subsubsection{Process supervisor}


\begin{itemize}
  \item \texttt{src/process/supervisor/types.ts} updated with discriminated spawn input and process local stdin contract.
  \item \texttt{src/process/supervisor/supervisor.ts} updated to use explicit \texttt{ptyCommand}.
  \item \texttt{src/process/supervisor/adapters/child.ts} and \texttt{src/process/supervisor/adapters/pty.ts} decoupled from agent types.
  \item \texttt{src/process/supervisor/registry.ts} idempotent finalize unchanged and retained.
\end{itemize}


\subsubsection{Exec and process integration}


\begin{itemize}
  \item \texttt{src/agents/bash-tools.exec-runtime.ts} updated to pass PTY command explicitly and keep fallback path.
  \item \texttt{src/agents/bash-tools.process.ts} updated to cancel via supervisor with real process-tree fallback termination.
  \item \texttt{src/agents/bash-tools.shared.ts} removed direct kill helper path.
\end{itemize}


\subsubsection{CLI reliability}


\begin{itemize}
  \item \texttt{src/agents/cli-watchdog-defaults.ts} added as shared baseline.
  \item \texttt{src/agents/cli-backends.ts} and \texttt{src/agents/cli-runner/reliability.ts} now consume same defaults.
\end{itemize}


\subsection{7. Validation run in this pass}


Unit tests:

\begin{itemize}
  \item \texttt{pnpm vitest src/process/supervisor/registry.test.ts}
  \item \texttt{pnpm vitest src/process/supervisor/supervisor.test.ts}
  \item \texttt{pnpm vitest src/process/supervisor/supervisor.pty-command.test.ts}
  \item \texttt{pnpm vitest src/process/supervisor/adapters/child.test.ts}
  \item \texttt{pnpm vitest src/agents/cli-backends.test.ts}
  \item \texttt{pnpm vitest src/agents/bash-tools.exec.pty-cleanup.test.ts}
  \item \texttt{pnpm vitest src/agents/bash-tools.process.poll-timeout.test.ts}
  \item \texttt{pnpm vitest src/agents/bash-tools.process.supervisor.test.ts}
  \item \texttt{pnpm vitest src/process/exec.test.ts}
\end{itemize}


E2E targets:

\begin{itemize}
  \item \texttt{pnpm vitest src/agents/cli-runner.test.ts}
  \item \texttt{pnpm vitest run src/agents/bash-tools.exec.pty-fallback.test.ts src/agents/bash-tools.exec.background-abort.test.ts src/agents/bash-tools.process.send-keys.test.ts}
\end{itemize}


Typecheck note:

\begin{itemize}
  \item Use \texttt{pnpm build} (and \texttt{pnpm check} for full lint/docs gate) in this repo. Older notes that mention \texttt{pnpm tsgo} are obsolete.
\end{itemize}


\subsection{8. Operational guarantees preserved}


\begin{itemize}
  \item Exec env hardening behavior is unchanged.
  \item Approval and allowlist flow is unchanged.
  \item Output sanitization and output caps are unchanged.
  \item PTY adapter still guarantees wait settlement on forced kill and listener disposal.
\end{itemize}


\subsection{9. Definition of done}


\begin{enumerate}
  \item Supervisor is lifecycle owner for managed runs.
  \item PTY spawn uses explicit command contract with no argv reconstruction.
  \item Process layer has no type dependency on agent layer for supervisor stdin contracts.
  \item Watchdog defaults are single source.
  \item Targeted unit and e2e tests remain green.
  \item Restart durability boundary is explicitly documented or fully implemented.
\end{enumerate}


\subsection{10. Summary}


The branch now has a coherent and safer supervision shape:

\begin{itemize}
  \item explicit PTY contract
  \item cleaner process layering
  \item supervisor driven cancellation path for process operations
  \item real fallback termination when supervisor lookup misses
  \item process-tree cancellation for child-run default kill paths
  \item unified watchdog defaults
  \item explicit in-memory restart boundary (no orphan reconciliation across restart in this pass)
\end{itemize}



\clearpage

% === 源文件: experiments/plans/session-binding-channel-agnostic.md ===
\section{Session Binding Channel Agnostic Plan}


\subsection{Overview}


This document defines the long term channel agnostic session binding model and the concrete scope for the next implementation iteration.

Goal:

\begin{itemize}
  \item make subagent bound session routing a core capability
  \item keep channel specific behavior in adapters
  \item avoid regressions in normal Discord behavior
\end{itemize}


\subsection{Why this exists}


Current behavior mixes:

\begin{itemize}
  \item completion content policy
  \item destination routing policy
  \item Discord specific details
\end{itemize}


This caused edge cases such as:

\begin{itemize}
  \item duplicate main and thread delivery under concurrent runs
  \item stale token usage on reused binding managers
  \item missing activity accounting for webhook sends
\end{itemize}


\subsection{Iteration 1 scope}


This iteration is intentionally limited.

\subsubsection{1. Add channel agnostic core interfaces}


Add core types and service interfaces for bindings and routing.

Proposed core types:

\begin{lstlisting}[language=typescript]
export type BindingTargetKind = "subagent" | "session";
export type BindingStatus = "active" | "ending" | "ended";

export type ConversationRef = {
  channel: string;
  accountId: string;
  conversationId: string;
  parentConversationId?: string;
};

export type SessionBindingRecord = {
  bindingId: string;
  targetSessionKey: string;
  targetKind: BindingTargetKind;
  conversation: ConversationRef;
  status: BindingStatus;
  boundAt: number;
  expiresAt?: number;
  metadata?: Record<string, unknown>;
};
\end{lstlisting}


Core service contract:

\begin{lstlisting}[language=typescript]
export interface SessionBindingService {
  bind(input: {
    targetSessionKey: string;
    targetKind: BindingTargetKind;
    conversation: ConversationRef;
    metadata?: Record<string, unknown>;
    ttlMs?: number;
  }): Promise<SessionBindingRecord>;

  listBySession(targetSessionKey: string): SessionBindingRecord[];
  resolveByConversation(ref: ConversationRef): SessionBindingRecord | null;
  touch(bindingId: string, at?: number): void;
  unbind(input: {
    bindingId?: string;
    targetSessionKey?: string;
    reason: string;
  }): Promise<SessionBindingRecord[]>;
}
\end{lstlisting}


\subsubsection{2. Add one core delivery router for subagent completions}


Add a single destination resolution path for completion events.

Router contract:

\begin{lstlisting}[language=typescript]
export interface BoundDeliveryRouter {
  resolveDestination(input: {
    eventKind: "task_completion";
    targetSessionKey: string;
    requester?: ConversationRef;
    failClosed: boolean;
  }): {
    binding: SessionBindingRecord | null;
    mode: "bound" | "fallback";
    reason: string;
  };
}
\end{lstlisting}


For this iteration:

\begin{itemize}
  \item only \texttt{task\_completion} is routed through this new path
  \item existing paths for other event kinds remain as-is
\end{itemize}


\subsubsection{3. Keep Discord as adapter}


Discord remains the first adapter implementation.

Adapter responsibilities:

\begin{itemize}
  \item create/reuse thread conversations
  \item send bound messages via webhook or channel send
  \item validate thread state (archived/deleted)
  \item map adapter metadata (webhook identity, thread ids)
\end{itemize}


\subsubsection{4. Fix currently known correctness issues}


Required in this iteration:

\begin{itemize}
  \item refresh token usage when reusing existing thread binding manager
  \item record outbound activity for webhook based Discord sends
  \item stop implicit main channel fallback when a bound thread destination is selected for session mode completion
\end{itemize}


\subsubsection{5. Preserve current runtime safety defaults}


No behavior change for users with thread bound spawn disabled.

Defaults stay:

\begin{itemize}
  \item \texttt{channels.discord.threadBindings.spawnSubagentSessions = false}
\end{itemize}


Result:

\begin{itemize}
  \item normal Discord users stay on current behavior
  \item new core path affects only bound session completion routing where enabled
\end{itemize}


\subsection{Not in iteration 1}


Explicitly deferred:

\begin{itemize}
  \item ACP binding targets (\texttt{targetKind: "acp"})
  \item new channel adapters beyond Discord
  \item global replacement of all delivery paths (\texttt{spawn\_ack}, future \texttt{subagent\_message})
  \item protocol level changes
  \item store migration/versioning redesign for all binding persistence
\end{itemize}


Notes on ACP:

\begin{itemize}
  \item interface design keeps room for ACP
  \item ACP implementation is not started in this iteration
\end{itemize}


\subsection{Routing invariants}


These invariants are mandatory for iteration 1.

\begin{itemize}
  \item destination selection and content generation are separate steps
  \item if session mode completion resolves to an active bound destination, delivery must target that destination
  \item no hidden reroute from bound destination to main channel
  \item fallback behavior must be explicit and observable
\end{itemize}


\subsection{Compatibility and rollout}


Compatibility target:

\begin{itemize}
  \item no regression for users with thread bound spawning off
  \item no change to non-Discord channels in this iteration
\end{itemize}


Rollout:

\begin{enumerate}
  \item Land interfaces and router behind current feature gates.
  \item Route Discord completion mode bound deliveries through router.
  \item Keep legacy path for non-bound flows.
  \item Verify with targeted tests and canary runtime logs.
\end{enumerate}


\subsection{Tests required in iteration 1}


Unit and integration coverage required:

\begin{itemize}
  \item manager token rotation uses latest token after manager reuse
  \item webhook sends update channel activity timestamps
  \item two active bound sessions in same requester channel do not duplicate to main channel
  \item completion for bound session mode run resolves to thread destination only
  \item disabled spawn flag keeps legacy behavior unchanged
\end{itemize}


\subsection{Proposed implementation files}


Core:

\begin{itemize}
  \item \texttt{src/infra/outbound/session-binding-service.ts} (new)
  \item \texttt{src/infra/outbound/bound-delivery-router.ts} (new)
  \item \texttt{src/agents/subagent-announce.ts} (completion destination resolution integration)
\end{itemize}


Discord adapter and runtime:

\begin{itemize}
  \item \texttt{src/discord/monitor/thread-bindings.manager.ts}
  \item \texttt{src/discord/monitor/reply-delivery.ts}
  \item \texttt{src/discord/send.outbound.ts}
\end{itemize}


Tests:

\begin{itemize}
  \item \texttt{src/discord/monitor/provider*.test.ts}
  \item \texttt{src/discord/monitor/reply-delivery.test.ts}
  \item \texttt{src/agents/subagent-announce.format.test.ts}
\end{itemize}


\subsection{Done criteria for iteration 1}


\begin{itemize}
  \item core interfaces exist and are wired for completion routing
  \item correctness fixes above are merged with tests
  \item no main and thread duplicate completion delivery in session mode bound runs
  \item no behavior change for disabled bound spawn deployments
  \item ACP remains explicitly deferred
\end{itemize}



\clearpage

